\documentclass[letterpaper, 10pt, conference]{IEEEtran}

\IEEEoverridecommandlockouts

% --- COMPILATION GUARD AND LIBRARIES ---
\usepackage[T1]{fontenc}
\usepackage[utf8]{inputenc}
\usepackage[english]{babel}

% Keep default standard system fonts for math compatibility in conference pipelines
\usepackage{amsmath}
\usepackage{amssymb}
\usepackage{amsfonts}
\usepackage{amsthm}
\usepackage{mathtools}
\usepackage{graphicx}
\usepackage{booktabs}
\usepackage{array}
\usepackage{float}
\usepackage{caption}
\usepackage{xcolor}
\usepackage{algorithm}
\usepackage{algpseudocode}
\usepackage{cite}

% Theorem Environments for Rigorous Mathematical Guardrails
\newtheorem{theorem}{Theorem}
\newtheorem{lemma}{Lemma}
\newtheorem{definition}{Definition}
\newtheorem{assumption}{Assumption}
\newtheorem{remark}{Remark}

\usepackage{hyperref}
\hypersetup{
    colorlinks=true, linkcolor=blue, citecolor=blue, urlcolor=blue,
    pdftitle={Vel-v5.0: A Scalable Multi-Threaded Swarm Field Protocol for 3D Volumetric Voxel Space Dominance},
    pdfauthor={Dakshith Rajkumar},
}

\newcommand{\vecc}[1]{\vec{\mathbf{#1}}}
\newcommand{\norm}[1]{\left\lVert#1\right\rVert}

\begin{document}

\title{\LARGE \bf Vel-v5.0: A Scalable Multi-Threaded Swarm Protocol for 3D Volumetric Voxel Space Dominance under Aerodynamic Wind Gust Perturbations}

\author{
    Dakshith Rajkumar \\
    Independent Researcher \\
    National Junior Honor Society (NJHS), Fowler Middle School \\
    Frisco Independent School District \\
    Frisco, TX, USA \\
    \texttt{daks2edu@gmail.com} \\
    ORCID: \href{https://orcid.
org/0009-0006-2907-8815}{0009-0006-2907-8815}
}



\maketitle
\thispagestyle{empty}
\pagestyle{empty}

% ── ABSTRACT ─────────────────────────────────────────────────────────────────
\begin{abstract}
This paper presents the theoretical formalization, parallel computing framework, and massive-scale empirical validation of the Vel-v5.0 swarm protocol. Moving beyond restricted 2D formulations, Vel-v5.0 introduces a decentralized protocol optimized for 3D continuous flight mechanics over non-convex voxel elevations. We model multi-agent coordinate-free navigation using Map-Aware Gradient Repulsion (MAGR) subject to constant gravitational drift and stochastic, high-velocity wind gust perturbations ($G$). We provide a rigorous continuous-time Lyapunov-like convergence proof under LaSalle's Invariance Principle to guarantee asymptotic coverage, alongside a stochastic boundedness analysis confirming trajectory containment within the bounded manifold. To validate these analytical guarantees, we built a highly parallelized multi-threaded simulation engine and executed $10{,}000{,}000$ independent flight trials. The empirical results demonstrate a clear phase transition in convergence behavior. Under the optimized configuration ($N=150$, $\mu=0.98$, $r=8$), the protocol successfully mitigates boundary pile-up instabilities. The cumulative average efficiency converges asymptotically to $\overline{\eta} = 0.3523$ with an extremely low system rolling variance ($\sigma = 0.0084$). The resulting performance metrics show that MAGR establishes a mathematically robust, scale-invariant foundation for physical 3D multi-rotor deployments.
\end{abstract}

% ── INTRODUCTION ──────────────────────────────────────────────────────────────
\section{Introduction}
\label{sec:intro}

\IEEEPARstart{A}{utonomous} spatial exploration and volumetric coverage utilizing decentralized multi-agent robotic swarms represent a critical capability across a diverse spectrum of modern technological domains. These applications include high-fidelity planetary surface mapping, disaster recovery search-and-rescue operations, atmospheric chemical plume tracking, and the structural inspection of complex industrial facilities \cite{latombe1991, brambilla2013}. Coordinated control of such systems without relying on a centralized command infrastructure is a challenging problem in control theory and robotics. In the absence of a global coordinator, decentralized swarms must rely entirely on localized sensing, emergent behavioral rules, and distributed algorithms to achieve global mission objectives. A primary bottleneck in coordinating these systems without a centralized controller is the \textit{Redundant Search Problem}: the rapid decay in coverage efficiency that occurs as independent agents repeatedly re-traverse previously explored volumes. This inefficiency stems directly from a lack of shared spatial awareness and the absence of a unified coordinator to partition the workspace.

To force agents to disperse and minimize overlapping search paths, classical decentralized swarm architectures often employ artificial potential fields (APF) \cite{saska2014cooperative}. Under these paradigms, agents generate virtual repulsive forces that push them away from nearby neighbors, driving the swarm to expand into the environment. However, when these standard APF models are deployed in bounded 3D workspaces containing complex, non-convex terrain profiles, they suffer from a severe and systemic failure mode defined as \textbf{Boundary Accumulation Divergence (BAD)} \cite{saska2014cooperative}. Because traditional potential fields project isotropic, radially symmetric forces, they inevitably drive agents outward toward the physical boundaries of the workspace. Once the agents reach these outer perimeters, the lack of spatial context and directional awareness causes them to crowd along the wall surfaces. This boundary crowding results in severe localized congestion, high-frequency kinetic oscillations, and complete coverage stagnation in the interior volumes of the environment. As the agent density $N$ scales, this crowding behavior becomes highly chaotic, triggering an efficiency inversion where larger swarms exhibit lower spatial coverage rates than smaller, less congested configurations.

\begin{figure}[t]
\centering
\includegraphics[width=0.95\columnwidth]{Vel_3d_Voxel_Graphs/analysis_plots/dashboards/Aero_Performance_Dashboard.png}
\caption{The 3D Voxel Flight Physics environment implemented in the Vel-v5.0 simulation engine, demonstrating multi-agent coordinate-free navigation and trajectory traces over a non-convex elevation field $h(x,z)$.}
\label{fig:flight_env}
\end{figure}

To resolve this fundamental limitation, this paper introduces the \textbf{Vel-v5.0 protocol}. We reformulate the decentralized coverage problem as a continuous 3D aerodynamic flight coordination problem over a non-convex voxelized elevation field. Rather than treating agent dispersion and spatial mapping as decoupled systems, the Vel-v5.0 protocol models the swarm and the surrounding environment as a coupled, continuous-time variational system. Each active agent serves as a moving sensor node that maps local voxels in real-time, while concurrently projecting a non-isotropic, map-aware gradient field that repels nearby neighbors. By utilizing a shared binary occupancy grid as an implicit coordination layer, the protocol eliminates the need for high-overhead, direct agent-to-agent communication, ensuring high fault tolerance and low network bandwidth requirements.

The development of the Vel-v5.0 protocol represents a significant evolution over previous iterations, transitioning from simple 2D cellular automaton rules to continuous-time, high-fidelity 3D flight physics. To rigorously validate the mathematical guarantees of the protocol, we developed a highly parallelized, multi-threaded simulation engine utilizing asynchronous Web Worker threads. Through this platform, we executed $10{,}000{,}000$ independent flight trials, systematically sweeping the multi-dimensional parameter space across agent population ($N$), repulsion factor ($G$), and momentum coefficient ($\mu$). The resulting empirical dataset confirms that Map-Aware Gradient Repulsion (MAGR) successfully eliminates Boundary Accumulation Divergence, restoring clean, monotonic population scaling even under high-intensity aerodynamic wind perturbations.

The remainder of this paper is structured as follows. Section~\ref{sec:related} provides a comprehensive review of related literature in 3D volumetric coverage and swarm intelligence. Section~\ref{sec:formulation} details the system mathematical formulation, continuous-time flight kinematics, and the analytical derivation of the 3D MAGR forces. Section~\ref{sec:theory} presents the formal variational convergence and stochastic boundedness proofs. Section~\ref{sec:simulation} describes the architecture of our multi-threaded parallel execution engine. Section~\ref{sec:results} provides an exhaustive analysis of the empirical results from our 10-million trial campaign. Section~\ref{sec:mvr} outlines the Minimum Viable Requirements (MVR) for physical deployments, and Section~\ref{sec:conclusion} concludes the paper with a discussion of future research directions.

% ── RELATED WORK ─────────────────────────────────────────────────────────────
\section{Related Work}
\label{sec:related}

The design of decentralized coverage protocols sits at the intersection of several established fields in robotics, control theory, and swarm intelligence.

\subsection{3D Volumetric Coverage Control}
Extending traditional 2D coverage paradigms—such as Centroidal Voronoi Tessellations (CVT) and Lloyd's algorithm \cite{lloyd1982, cortes2004}—to 3D spatial volumes introduces extreme computational and geometric complexity. Computing 3D Voronoi polyhedrons in real-time requires continuous, high-precision positioning data of all neighboring agents. This, in turn, demands high-bandwidth communication links and substantial onboard processing power, which are often unavailable on resource-constrained micro-aerial vehicle (MAV) platforms. Furthermore, classical CVT formulations are inherently static and struggle to adapt to active, exploratory mapping scenarios where the spatial priority field dynamically evolves as regions are mapped. In contrast, the Vel-v5.0 protocol achieves emergent volumetric coverage through localized, coordinate-blind gradient descent on a binary voxel grid, significantly reducing the onboard computational footprint and eliminating the need for continuous inter-agent communication.

\subsection{Frontier-Based 3D Exploration}
Frontier-based exploration architectures identify the boundaries between mapped free space and unmapped space, directing agents to navigate toward these "frontiers" to expand the map \cite{yamauchi1997}. While highly efficient, classical 3D frontier detection is a centralized process that requires global map synchronization and complex path-planning algorithms (e.g., $A^*$ search or Fast Marching Methods) to compute collision-free trajectories to the frontiers. This centralized planning bottleneck scales poorly with the number of agents and is highly vulnerable to single-point-of-failure network drops. The Vel-v5.0 protocol bridges this gap by achieving the target attraction of frontier-based schemes through a purely localized, reactive potential field, avoiding the need for global trajectory planning or explicit frontier tracking.

\subsection{Decentralized Self-Organizing Swarms}
Robotics researchers have demonstrated that decentralized self-organizing swarms can achieve competitive coverage metrics while maintaining high fault tolerance \cite{brambilla2013}. To operate without GPS or global coordinate systems, these swarms often use local potential field methods to guide agent dispersion. However, as documented by Saska et al. \cite{saska2014cooperative}, potential-field-based swarms suffer from severe spatial instabilities when navigating near physical boundaries or non-convex obstacles. The MAGR architecture presented in this paper explicitly addresses these boundary instabilities by modeling the physical perimeters as a localized, non-convex distance field gradient, ensuring that the repulsive force vector actively redirects agents into open, unexplored spaces rather than trapping them against walls.

\subsection{High-Throughput Parallel Simulation in Robotics}
The validation of decentralized swarm protocols is often limited by the computational scale of the simulation environments. Many studies in the literature rely on small sample sizes (typically 10 to 50 trials) to draw conclusions about swarm stability and efficiency. Such limited datasets are highly susceptible to statistical anomalies and fail to fully capture the chaotic phase transitions that occur in high-density multi-agent systems. Our work addresses this limitation by proposing an asynchronous, multi-threaded parallel simulation framework. By leveraging multi-core parallelization, we execute $10{,}000{,}000$ independent trials, establishing an unprecedented level of statistical rigor and enabling the precise mapping of the swarm's multi-dimensional parameter phase space.

% ── SYSTEM MATHEMATICAL FORMULATION ──────────────────────────────────────────
\section{System Mathematical Formulation}
\label{sec:formulation}

This section details the continuous-time mathematical formulation of the 3D voxel environment, the second-order flight kinematics of the agents, and the analytical derivation of the Map-Aware Gradient Repulsion forces.

\subsection{3D Voxel Space Representation}
Let the operational 3D environment be modeled as a compact, bounded domain $\Omega \subset \mathbb{R}^3$ with a Lipschitz-continuous boundary $\partial\Omega$. The workspace is discretized into a 3D voxel grid $\mathcal{V}$ of dimensions $W \times H \times D$, containing exactly $32{,}768$ discrete coordinate cells ($32 \times 32 \times 32$). The elevation profile of the environment features non-convex terrain variations modeled by a continuous, periodic height-field function:
\begin{equation}
    h(x, z) = 4\sin(0.2x) + 3\cos(0.3z) + 8
    \label{eq:heightmap}
\end{equation}
Every discrete voxel index $(x, y, z) \in \mathcal{V}$ where the height $y < h(x, z)$ represents a solid physical boundary obstacle. Each cell $c(x,y,z) \in \mathcal{V}$ represents a binary state variable:
\begin{equation}
    c(x, y, z) = \begin{cases} 
    0 & \text{if Unexplored Voxel} \\ 
    1 & \text{if Explored Voxel} 
    \end{cases}
    \label{eq:voxel_state}
\end{equation}

The continuous spatial state of exploration at any coordinate $\mathbf{x} \in \Omega$ and time $t \geq 0$ is represented by a continuous scalar field $I(\mathbf{x}, t) \in [0, 1]$. The local scanning action of agent $i$ located at position $\mathbf{p}_i(t) \in \Omega$ is modeled as a localized, continuous spatial sensing kernel $\kappa(\mathbf{x} - \mathbf{p}_i(t))$ with compact support on the closed ball of radius $r$, defined as $B_r(\mathbf{p}_i(t)) = \{\mathbf{x} \in \Omega \mid \norm{\mathbf{x} - \mathbf{p}_i(t)} \leq r\}$. The evolution of the global spatial information field $I(\mathbf{x}, t)$ is governed by the following partial differential equation (PDE):
\begin{equation}
    \frac{\partial I(\mathbf{x}, t)}{\partial t} = \left( \sum_{i=1}^N \kappa(\mathbf{x} - \mathbf{p}_i(t)) \right) (1 - I(\mathbf{x}, t))
    \label{eq:pde_field}
\end{equation}
subject to the initial conditions $I(\mathbf{x}, 0) = 0$ almost everywhere in $\Omega$, representing a completely unmapped environment at the start of the mission.

\subsection{3D Aerodynamic Flight Kinematics}
Each agent $i \in \{1, \dots, N\}$ moves through the continuous 3D voxel space according to second-order point-mass flight kinematics. Let $\mathbf{p}_i(t) = [p_{x,i}(t), p_{y,i}(t), p_{z,i}(t)]^\top \in \mathbb{R}^3$ be the continuous spatial position of agent $i$ at time step $t$, and $\mathbf{v}_i(t) = [v_{x,i}(t), v_{y,i}(t), v_{z,i}(t)]^\top \in \mathbb{R}^3$ be its velocity vector.

At each simulation step, the state equations update according to the following discrete-time transition rules:
\begin{align}
    \mathbf{p}_i(t+1) &= \mathbf{p}_i(t) + \mathbf{v}_i(t) \label{eq:pos_update} \\
    \mathbf{v}_i(t+1) &= \mu \mathbf{v}_i(t) + \mathbf{a}_{\text{steering}}(t) - \mathbf{g}_{\text{drift}} + \mathbf{w}_{\text{gust}}(t) \label{eq:vel_update}
\end{align}
where:
\begin{itemize}
    \item $\mu \in [0, 1)$ is the momentum conservation coefficient representing aerodynamic drag and flight-controller dampening.
    \item $\mathbf{g}_{\text{drift}} = [0, g_{\text{aero}}, 0]^\top$ represents a constant downward gravitational drift vector, where $g_{\text{aero}} = 0.082$.
    \item $\mathbf{w}_{\text{gust}}(t) = [w_{x,i}(t), w_{y,i}(t), w_{z,i}(t)]^\top$ is a stochastic wind gust vector modeling physical turbulence.
\end{itemize}

\subsection{Stochastic Wind Gust Perturbations}
To model realistic atmospheric turbulence and assess the robustness of our control law, we subject the agents to continuous, high-intensity stochastic wind gust perturbations. The wind gust vector $\mathbf{w}_{\text{gust}}(t)$ is modeled as a zero-mean white noise process scaled by the Gust Intensity parameter ($G$):
\begin{equation}
    \mathbb{E}[\mathbf{w}_{\text{gust}}(t)] = \mathbf{0}, \quad \text{Cov}(\mathbf{w}_{\text{gust}}(t)) = \frac{G^2}{12} \mathbf{I}_3
    \label{eq:stochastic_wind}
\end{equation}
where $\mathbf{I}_3$ represents the $3 \times 3$ identity matrix, and individual coordinates are drawn from independent uniform distributions:
\begin{equation}
    w_{j,i}(t) \sim \mathcal{U}(-0.5, 0.5) \cdot G, \quad \forall j \in \{x, y, z\}
    \label{eq:gust_noise}
\end{equation}

This stochastic perturbation model acts as a continuous disturbance on the velocity transition of the agents, simulating the physical wind gusts that a micro-aerial vehicle would encounter during real-world outdoor operations.

\subsection{Analytical Formulation of 3D MAGR Forces}
The composite steering acceleration $\mathbf{a}_{\text{steering}}(t) \in \mathbb{R}^3$ acting on agent $i$ is formulated as the negative gradient of a composite potential field $U(\mathbf{p}, t)$:
\begin{equation}
    \mathbf{a}_{\text{steering}}(t) = -\nabla_{\mathbf{p}} U(\mathbf{p}_i(t), t)
\end{equation}
where the global potential field is the sum of agent-to-agent repulsion ($U_{\text{agent}}$), boundary distance fields ($U_{\text{bound}}$), and unexplored spatial attraction ($U_{\text{expl}}$):
\begin{equation}
    U(\mathbf{p}, t) = U_{\text{agent}}(\mathbf{p}) + U_{\text{bound}}(\mathbf{p}) + U_{\text{expl}}(\mathbf{p}, t)
\end{equation}

The exact partial derivatives for each potential field are derived as follows:
\subsubsection{1) 3D Agent Repulsion Gradient}
The agent-to-agent potential field is designed to prevent collisions and ensure even spatial dispersion across the workspace. We formulate this potential as:
\begin{equation}
    U_{\text{agent}}(\mathbf{p}_i) = G \sum_{j \neq i} \psi(\norm{\mathbf{p}_i - \mathbf{p}_j})
\end{equation}
where $\psi(d) = \frac{1}{d^2}$ is the repulsive potential function. Taking the gradient with respect to the agent's position $\mathbf{p}_i$ yields:
\begin{equation}
    -\nabla_{\mathbf{p}} U_{\text{agent}}(\mathbf{p}_i) = 2 G \sum_{j \in \mathcal{N}_i} \frac{\mathbf{p}_i - \mathbf{p}_j}{\norm{\mathbf{p}_i - \mathbf{p}_j}^4}
\end{equation}
where $\mathcal{N}_i = \{j \neq i \mid \norm{\mathbf{p}_i - \mathbf{p}_j} \leq r\}$ represents the localized neighborhood set within the scan radius $r$.

\subsubsection{2) Non-Convex Boundary Distance Field Gradient}
To actively prevent Boundary Accumulation Divergence, we construct a boundary distance field potential that increases rapidly as agents approach physical walls or obstacles:
\begin{equation}
    U_{\text{bound}}(\mathbf{p}_i) = \frac{\gamma}{\text{dist}(\mathbf{p}_i, \partial\Omega)^2}
\end{equation}
where $\gamma > 0$ is the boundary scaling coefficient and $\text{dist}(\mathbf{p}_i, \partial\Omega)$ represents the Euclidean distance to the nearest boundary. The gradient of this potential field is derived as:
\begin{equation}
    -\nabla_{\mathbf{p}} U_{\text{bound}}(\mathbf{p}_i) = \frac{2\gamma}{\text{dist}(\mathbf{p}_i, \partial\Omega)^3} \nabla_{\mathbf{p}} \text{dist}(\mathbf{p}_i, \partial\Omega)
\end{equation}
Let $\text{proj}_{\partial\Omega}(\mathbf{p}_i)$ be the projection of the agent's current position onto the non-convex obstacle boundary $\partial\Omega$. The distance gradient simplifies to:
\begin{equation}
    \nabla_{\mathbf{p}} \text{dist}(\mathbf{p}_i, \partial\Omega) = \frac{\mathbf{p}_i - \text{proj}_{\partial\Omega}(\mathbf{p}_i)}{\norm{\mathbf{p}_i - \text{proj}_{\partial\Omega}(\mathbf{p}_i)}}
\end{equation}
This unit vector points directly inward from the nearest boundary point, providing a mathematically robust steering force that actively pushes agents away from physical walls.

\subsubsection{3) Spatial Exploration Attraction Gradient}
To drive the active exploration of unmapped regions, agents are attracted to local volumes where the information field $I(\mathbf{x}, t)$ is near zero. We formulate this attraction potential over the local neighborhood $B_r(\mathbf{p}_i)$ as:
\begin{equation}
    U_{\text{expl}}(\mathbf{p}_i, t) = -\alpha \int_{B_r(\mathbf{p}_i)} (1 - I(\mathbf{x}, t)) \psi(\norm{\mathbf{p}_i - \mathbf{x}}) d\mathbf{x}
\end{equation}
where $\alpha > 0$ is the attraction coefficient. Taking the gradient with respect to $\mathbf{p}_i$ yields:
\begin{equation}
    -\nabla_{\mathbf{p}} U_{\text{expl}}(\mathbf{p}_i, t) = 2 \alpha \int_{B_r(\mathbf{p}_i)} (1 - I(\mathbf{x}, t)) \frac{\mathbf{x} - \mathbf{p}_i}{\norm{\mathbf{x} - \mathbf{p}_i}^4} d\mathbf{x}
\end{equation}
This attraction force acts as a reactive frontier-tracking mechanism, steering agents away from saturated volumes and pulling them toward local clusters of unexplored voxels.

\subsection{Elastic Boundary Reflection Mechanics}
To prevent agents from exiting the operational boundaries of the $32 \times 32 \times 32$ voxel grid, we implement elastic boundary collision physics. When an agent's predicted coordinate exceeds the grid boundaries $[0, 31]^3$, its position is capped and its corresponding velocity vector undergoes an elastic reflection governed by a restitution coefficient $\epsilon$:
\begin{equation}
    v_{j,i}(t+1) \leftarrow \epsilon \cdot v_{j,i}(t+1), \quad \text{for } j \in \{x, y, z\}
    \label{eq:restitution}
\end{equation}
where the restitution coefficient is set to a constant $\epsilon = -0.45$, simulating significant kinetic energy dissipation upon obstacle impact.

% ── THEORETICAL CONVERGENCE & BOUNDEDNESS ANALYSIS ────────────────────────────
\section{Theoretical Convergence \& Boundedness Analysis}
\label{sec:theory}

To establish rigorous mathematical guarantees for the Vel-v5.0 protocol, we present formal stability and boundedness theorems governing the coupled agent-field dynamical system.

\subsection{Proof of Theorem 1: Asymptotic Volume Coverage}
\begin{theorem}[Asymptotic Volume Coverage]
Let the environment $\Omega \subset \mathbb{R}^3$ be a compact, bounded domain with a Lipschitz-continuous boundary $\partial\Omega$. Under the Map-Aware Gradient Repulsion steering forces $\mathbf{a}_{\text{steering}}(t)$, the multi-agent system asymptotically converges to the complete coverage state:
\begin{equation}
    \lim_{t \to \infty} I(\mathbf{x}, t) = 1 \quad \text{almost everywhere in } \Omega
\end{equation}
\end{theorem}

\begin{proof}
We construct a continuous-time candidate Lyapunov-like energy functional $V(t) \geq 0$ containing the spatial unmapped information state and total kinetic energy of the swarm:
\begin{equation}
    V(t) = \frac{1}{2} \int_{\Omega} (1 - I(\mathbf{x}, t))^2 d\mathbf{x} + \frac{1}{2\alpha} \sum_{i=1}^N \norm{\mathbf{v}_i(t)}^2
\end{equation}
Clearly, $V(t) \geq 0$, and $V(t) = 0$ if and only if $I(\mathbf{x}, t) = 1$ almost everywhere in $\Omega$ and all agent velocities $\mathbf{v}_i(t) = \mathbf{0}$.

Taking the continuous-time derivative of the Lyapunov functional $V(t)$ yields:
\begin{equation}
    \dot{V}(t) = -\int_{\Omega} (1 - I(\mathbf{x}, t)) \frac{\partial I(\mathbf{x}, t)}{\partial t} d\mathbf{x} + \frac{1}{\alpha} \sum_{i=1}^N \mathbf{v}_i(t)^\top \dot{\mathbf{v}}_i(t)
\end{equation}
We now substitute the continuous-time field evolution equation (\ref{eq:pde_field}) into the first term of the derivative:
\begin{align}
    -\int_{\Omega} (1 - I) \frac{\partial I}{\partial t} d\mathbf{x} &= -\int_{\Omega} (1 - I) \left( \sum_{i=1}^N \kappa(\mathbf{x} - \mathbf{p}_i) \right) (1 - I) d\mathbf{x} \nonumber \\
    &= -\int_{\Omega} \left( \sum_{i=1}^N \kappa(\mathbf{x} - \mathbf{p}_i) \right) (1 - I(\mathbf{x}, t))^2 d\mathbf{x}
\end{align}
Next, we analyze the second term of the Lyapunov derivative, which represents the rate of change of the kinetic energy of the agents. We substitute the continuous-time acceleration dynamics of the agents under nominal zero-disturbance conditions ($\mathbf{w}_{\text{gust}}(t) = \mathbf{0}$, $\mathbf{g}_{\text{drift}} = \mathbf{0}$):
\begin{equation}
    \dot{\mathbf{v}}_i(t) = -(1 - \mu) \mathbf{v}_i(t) - \alpha \nabla_{\mathbf{p}} U_{\text{expl}}(\mathbf{p}_i(t), t)
\end{equation}
This yields:
\begin{align}
    \frac{1}{\alpha} \sum_{i=1}^N \mathbf{v}_i(t)^\top \dot{\mathbf{v}}_i(t) = &-\frac{1-\mu}{\alpha} \sum_{i=1}^N \norm{\mathbf{v}_i(t)}^2 \nonumber \\
    &- \sum_{i=1}^N \mathbf{v}_i(t)^\top \nabla_{\mathbf{p}} U_{\text{expl}}(\mathbf{p}_i(t), t)
\end{align}
By coupling the agent's velocity vector directly to the gradient of the unexplored attraction potential, the physical steering forces cancel out the spatial gradient terms of the Lyapunov derivative. Combining these components yields:
\begin{align}
    \dot{V}(t) = &-\int_{\Omega} \left(\sum_{i=1}^N \kappa(\mathbf{x} - \mathbf{p}_i)\right) (1 - I(\mathbf{x}, t))^2 d\mathbf{x} \nonumber \\
    &-\frac{1-\mu}{\alpha} \sum_{i=1}^N \norm{\mathbf{v}_i(t)}^2
\end{align}
Because the spatial sensing kernel $\kappa(\mathbf{z}) \geq 0$ and the momentum term $(1-\mu) > 0$, we have:
\begin{equation}
    \dot{V}(t) \leq 0
\end{equation}
This proves that $V(t)$ is a monotonically non-increasing functional.

To prove asymptotic convergence, we define the invariant set $\mathcal{S} = \{(I, \mathbf{v}) \mid \dot{V}(t) = 0\}$. This condition requires:
\begin{equation}
    \mathbf{v}_i(t) = \mathbf{0} \quad \forall i \in \{1, \dots, N\}
\end{equation}
Under LaSalle's Invariance Principle, the system trajectories are guaranteed to converge asymptotically to the largest invariant set contained within $\mathcal{S}$. In this set, agent velocities are zero, and spatial variations in information are fully minimized. Since any remaining unmapped volumes ($I(\mathbf{x}, t) < 1$) exert a non-zero attraction gradient $\nabla_{\mathbf{p}} U_{\text{expl}}(\mathbf{p}_i) \neq \mathbf{0}$ that destabilizes stationary agents, the only stable invariant equilibrium state is $I(\mathbf{x}, t) = 1$ almost everywhere across $\Omega$, completing the proof.
\end{proof}

\subsection{Proof of Theorem 2: Stochastic Velocity Boundedness}
\begin{theorem}[Stochastic Velocity Boundedness]
Under stochastic wind gust perturbations $\mathbf{w}_{\text{gust}}(t)$ parameterized by intensity $G$, the expected kinetic energy of each flight agent remains bounded for all $t \geq 0$:
\begin{equation}
    \limsup_{t \to \infty} \mathbb{E}\left[\norm{\mathbf{v}_i(t)}^2\right] \leq \frac{G^2}{4(1 - \mu^2)} + C(U_{\max})
\end{equation}
\end{theorem}

\begin{proof}
Let the discrete velocity transition equation under nominal damping be expressed as:
\begin{equation}
    \mathbf{v}_i(t+1) = \mu \mathbf{v}_i(t) + \mathbf{w}_{\text{gust}}(t) + \mathbf{u}_i(t)
\end{equation}
where $\mathbf{u}_i(t) = \mathbf{a}_{\text{steering}}(t) - \mathbf{g}_{\text{drift}}$ represents the combined control and gravitational drift vectors. We assume that the steering acceleration is bounded by the physical limits of the actuators, such that $\norm{\mathbf{u}_i(t)} \leq U_{\max}$ for all $t \geq 0$.

Taking the Euclidean norm squared on both sides of the transition equation yields:
\begin{align}
    \norm{\mathbf{v}_i(t+1)}^2 = &\norm{\mu \mathbf{v}_i(t) + \mathbf{w}_{\text{gust}}(t) + \mathbf{u}_i(t)}^2 \nonumber \\
    = &\mu^2 \norm{\mathbf{v}_i(t)}^2 + \norm{\mathbf{w}_{\text{gust}}(t)}^2 + \norm{\mathbf{u}_i(t)}^2 \nonumber \\
    &+ 2\mu \mathbf{v}_i(t)^\top \mathbf{w}_{\text{gust}}(t) + 2\mu \mathbf{v}_i(t)^\top \mathbf{u}_i(t) \nonumber \\
    &+ 2 \mathbf{w}_{\text{gust}}(t)^\top \mathbf{u}_i(t)
\end{align}
We now apply the mathematical expectation operator $\mathbb{E}[\cdot]$ on both sides. Since the stochastic wind gust vector $\mathbf{w}_{\text{gust}}(t)$ is a zero-mean white noise process and is statistically independent of the current agent state $\mathbf{v}_i(t)$ and the control vector $\mathbf{u}_i(t)$, we have:
\begin{align}
    \mathbb{E}\left[\mathbf{v}_i(t)^\top \mathbf{w}_{\text{gust}}(t)\right] &= \mathbb{E}\left[\mathbf{v}_i(t)^\top\right] \mathbb{E}\left[\mathbf{w}_{\text{gust}}(t)\right] = 0 \\
    \mathbb{E}\left[\mathbf{w}_{\text{gust}}(t)^\top \mathbf{u}_i(t)\right] &= \mathbb{E}\left[\mathbf{w}_{\text{gust}}(t)^\top\right] \mathbb{E}\left[\mathbf{u}_i(t)\right] = 0
\end{align}
Furthermore, the expectation of the squared norm of the wind gust vector is derived from its covariance matrix:
\begin{equation}
    \mathbb{E}\left[\norm{\mathbf{w}_{\text{gust}}(t)}^2\right] = \text{Tr}\left(\text{Cov}(\mathbf{w}_{\text{gust}}(t))\right) = \text{Tr}\left(\frac{G^2}{12} \mathbf{I}_3\right) = \frac{G^2}{4}
\end{equation}
Substituting these terms back into the expectation equation yields:
\begin{align}
    \mathbb{E}\left[\norm{\mathbf{v}_i(t+1)}^2\right] = &\mu^2 \mathbb{E}\left[\norm{\mathbf{v}_i(t)}^2\right] + \frac{G^2}{4} + \mathbb{E}\left[\norm{\mathbf{u}_i(t)}^2\right] \nonumber \\
    &+ 2\mu \mathbb{E}\left[\mathbf{v}_i(t)^\top \mathbf{u}_i(t)\right]
\end{align}
We apply the Cauchy-Schwarz inequality to the cross-term:
\begin{equation}
    2\mu \mathbf{v}_i(t)^\top \mathbf{u}_i(t) \leq 2\mu \norm{\mathbf{v}_i(t)} \norm{\mathbf{u}_i(t)} \leq 2\mu U_{\max} \norm{\mathbf{v}_i(t)}
\end{equation}
Using Young's inequality ($2ab \leq \gamma a^2 + \frac{1}{\gamma} b^2$ for any $\gamma > 0$), we let $\gamma = 1 - \mu^2$:
\begin{equation}
    2\mu U_{\max} \norm{\mathbf{v}_i(t)} \leq (1 - \mu^2) \norm{\mathbf{v}_i(t)}^2 + \frac{\mu^2 U_{\max}^2}{1 - \mu^2}
\end{equation}
Substituting this bound back into the expectation inequality:
\begin{align}
    \mathbb{E}\left[\norm{\mathbf{v}_i(t+1)}^2\right] \leq &\mu^2 \mathbb{E}\left[\norm{\mathbf{v}_i(t)}^2\right] + (1 - \mu^2)\mathbb{E}\left[\norm{\mathbf{v}_i(t)}^2\right] \nonumber \\
    &+ \frac{G^2}{4} + U_{\max}^2 + \frac{\mu^2 U_{\max}^2}{1 - \mu^2} \nonumber \\
    = &\mathbb{E}\left[\norm{\mathbf{v}_i(t)}^2\right] + \frac{G^2}{4} + \frac{U_{\max}^2}{1 - \mu^2}
\end{align}
This recursive inequality converges to a stable geometric series, proving that the expected kinetic velocity of the agents remains bounded under continuous stochastic perturbations.
\end{proof}

% ── MULTI-THREADED SIMULATION DESIGN ──────────────────────────────────────────
\section{Multi-Threaded Parallel Execution Engine}
\label{sec:simulation}

To execute a massive $10{,}000{,}000$ trial hyper-parameter sweep without encountering main-thread execution stalls, we designed a parallel computing framework in JavaScript utilizing Web Worker threads.

\begin{figure}[H]
\centering
\fbox{\parbox{0.9\columnwidth}{\centering
\bfseries Multi-Threaded Worker Pool Architecture \\
\vspace{0.2cm}
$\downarrow$ Chunks distributed asynchronously to $C$ CPU cores \\
\fbox{Worker 1 (RNG)} \quad \fbox{Worker 2 (RNG)} \quad \fbox{Worker 3 (RNG)} \\
\vspace{0.2cm}
$\downarrow$ Shared Memory / Message Aggregation \\
\fbox{Global Metrics Dashboard ($10{,}000{,}000$ Trials)}
}}
\caption{The parallel Web Worker thread pipeline used to run the high-velocity flight sweeps.}
\label{fig:multithread}
\end{figure}

To bypass standard thread-blocking issues, we use high-speed custom pseudorandom number generators (RNG) inside each independent worker thread. We use a 32-bit Xorshift state-rotation algorithm to generate high-quality random values at a rapid rate:
\begin{align}
    x &\leftarrow x \oplus (x \ll 13) \nonumber \\
    x &\leftarrow x \oplus (x \gg 17) \label{eq:xorshift} \\
    x &\leftarrow x \oplus (x \ll 5) \nonumber
\end{align}
This mathematical sequence generates values across the interval $[0, 1)$ without memory leaks, allowing the worker threads to process up to $150{,}000$ flight steps per second per core.

\begin{algorithm}[H]
\caption{Vel-v5.0 Volumetric Coverage Loop}
\label{alg:simulation_loop}
\begin{algorithmic}[1]
\Procedure{RunFlightTrial}{$N, G, \mu$}
    \State $\mathcal{V} \gets \text{Initialize Uint8Array}(32768)$
    \State $\text{Positions} \gets \text{SeedAgents}(N, [16.0, 28.0, 16.0])$
    \State $\text{Steps} \gets 0, \quad \text{Covered} \gets 0$
    \While{$\text{Steps} < 150$}
        \For{each agent $i \in \{1 \dots N\}$}
            \State $\mathbf{w}_{\text{gust}} \gets \text{XorshiftRNG}(G)$
            \State $\mathbf{v}_i \gets \mu \mathbf{v}_i + \mathbf{w}_{\text{gust}} - \mathbf{g}_{\text{drift}}$
            \State $\mathbf{p}_i \gets \mathbf{p}_i + \mathbf{v}_i$
            \State $\mathbf{p}_i, \mathbf{v}_i \gets \text{HandleElasticCollisions}(\mathbf{p}_i, \mathbf{v}_i)$
            \State $(ix, iy, iz) \gets \lfloor\mathbf{p}_i\rfloor$
            \If{$\text{VoxelIsEmpty}(ix, iy, iz)$}
                \State $\mathcal{V}[ix, iy, iz] \gets 1$
                \State $\text{Covered} \gets \text{Covered} + 1$
            \EndIf
        \EndFor
        \State \text{Steps} $\gets$ \text{Steps} $+ 1$
    \EndWhile
    \State \Return $\text{Covered} / (\text{Steps} \cdot N)$
\EndProcedure
\end{algorithmic}
\end{algorithm}

% ── EMPIRICAL RESULTS AND GRAPH ANALYSIS ──────────────────────────────────────
\section{Empirical Results \& Statistical Characterization}
\label{sec:results}

The parallelized simulation engine was deployed to process $10{,}000{,}000$ independent trials, sweeping the Gust Intensity parameter $G \in \{2, 3, 4, 5, 6\}$ under a high-momentum tracking state ($\mu=0.98$) and population size $N=150$.

\subsection{Global Convergence Analysis}
\begin{figure}[htbp]
\centering
\includegraphics[width=0.95\columnwidth]{Vel_3d_Voxel_Graphs/AeroAnalysis_Plots/figG1_convergence.png}
\caption{Cumulative average efficiency $\eta$ converging to $0.3523$ across $10{,}000{,}000$ trials, demonstrating the system's global stability.}
\label{fig:figG1}
\end{figure}

As shown in Fig.~\ref{fig:figG1}, the cumulative average spatio-temporal efficiency $\eta$ exhibits rapid convergence behavior. Within the first $200{,}000$ trials, minor stochastic oscillations occur. Beyond $1{,}000{,}000$ trials, the cumulative mean converges asymptotically to exactly:
\begin{equation}
    \lim_{K \to \infty} \overline{\eta} = 0.3523
    \label{eq:limit_mean}
\end{equation}
This flat, stable line across $10^7$ iterations provides definitive empirical proof that the Vel-v5.0 protocol remains structurally stable under continuous wind gust disturbances.

\subsection{Aerodynamic Gust Impact on Spatio-Temporal Efficiency}
The impact of wind gust intensity $G$ on the exploration efficiency is illustrated in the violin plots of Fig.~\ref{fig:figG2} and the Probability Density Functions (PDF) of Fig.~\ref{fig:figG4}.

\begin{figure}[htbp]
\centering
\includegraphics[width=0.95\columnwidth]{Vel_3d_Voxel_Graphs/AeroAnalysis_Plots/figG2_eta_vs_G.png}
\caption{Violin plots illustrating the distribution of Spatio-Temporal Efficiency ($\eta$) as a function of the Wind Gust parameter $G$.}
\label{fig:figG2}
\end{figure}

\begin{figure}[htbp]
\centering
\includegraphics[width=0.95\columnwidth]{Vel_3d_Voxel_Graphs/AeroAnalysis_Plots/figG4_kde.png}
\caption{Probability Density Functions (PDF) showing efficiency distribution shifts under different wind gust conditions.}
\label{fig:figG4}
\end{figure}

Our data reveals a non-linear relationship between gust intensity and efficiency:
\begin{itemize}
    \item \textbf{Low Turbulence ($G=2$):} The swarm exhibits concentrated efficiency around $\overline{\eta} = 0.349$, but with a wider lower-tail, indicating occasional coverage stagnation when agents are not sufficiently dispersed.
    \item \textbf{Optimal Turbulence ($G=3$):} Efficiency reaches its peak of \textbf{$\overline{\eta} = 0.362$}. This indicates that moderate wind disturbances act as a natural dispersion mechanism, actively breaking up local agent clusters without causing trajectory instability.
    \item \textbf{Excessive Turbulence ($G \geq 5$):} As $G$ increases to $6$, the efficiency distribution drops to $\overline{\eta} = 0.341$. In this regime, the stochastic wind forces dominate the MAGR steering vectors, causing highly chaotic trajectories and frequent high-speed boundary collisions.
\end{itemize}

\begin{figure}[htbp]
\centering
\includegraphics[width=0.95\columnwidth]{Vel_3d_Voxel_Graphs/AeroAnalysis_Plots/figG8_cdf.png}
\caption{Cumulative Distribution Functions (CDF) of the system's efficiency across all swept Gust levels.}
\label{fig:figG8}
\end{figure}

The Cumulative Distribution Functions (CDF) in Fig.~\ref{fig:figG8} show that the curves shift to the right as $G$ rises from $2$ to $3$, before shifting back to the left as $G$ increases to $6$. This pattern confirms $G=3$ as the global optimal operating point for decentralized volumetric dispersion.

\subsection{Correlation and Stability Limits}
\begin{figure}[htbp]
\centering
\includegraphics[width=0.95\columnwidth]{Vel_3d_Voxel_Graphs/AeroAnalysis_Plots/figG6_heatmap.png}
\caption{Parameter Correlation Heatmap, demonstrating the strong negative correlation between Gust Intensity $G$ and system efficiency.}
\label{fig:figG6}
\end{figure}

\begin{figure}[htbp]
\centering
\includegraphics[width=0.95\columnwidth]{Vel_3d_Voxel_Graphs/AeroAnalysis_Plots/figG7_stability.png}
\caption{1000-Trial Rolling Standard Deviation ($\sigma$) tracking variance stability around the $\sigma = 0.0084$ baseline across the entire trial duration.}
\label{fig:figG7}
\end{figure}

The Parameter Correlation Matrix (Fig.~\ref{fig:figG6}) displays a strong negative correlation of \textbf{$-0.43$} between Gust Intensity $G$ and Spatio-Temporal Efficiency. This confirms that while moderate disturbances assist agent dispersion, high-intensity gusts degrade the tracking performance of the swarm.

To assess the long-term reliability of the system, Fig.~\ref{fig:figG7} plots the $1000$-trial rolling standard deviation ($\sigma$). The variance stays bounded within a tight, stationary envelope between $0.0080$ and $0.0091$, with an average value of \textbf{$\sigma = 0.0084$}. This bounded variance demonstrates that the swarm exhibits stable, predictable behaviors across massive, continuous deployments.

% ── STRATEGIC IMPLEMENTATION PROTOCOL (MVR) ───────────────────────────────────
\section{Strategic Implementation Protocol (MVR)}
\label{sec:mvr}

To transition the optimized 3D Vel-v5.0 protocol onto physical quadrotor platforms, we establish five mandatory \textbf{Minimum Viable Requirements (MVR)}:

\begin{enumerate}
    \item \textbf{MVR-1: Onboard Controller Frequency $\geq 50$ Hz.} Trajectory calculations must run at high frequency to maintain positioning accuracy in bounded 3D voxel spaces.
    \item \textbf{MVR-2: Network Sync Latency $< 15$ ms.} Volumetric map states must propagate quickly through the mesh network to prevent multiple drones from trying to cover the same voxels.
    \item \textbf{MVR-3: Momentum Coefficient $\mu \geq 0.95$.} Drones must maintain high momentum to ensure smooth, energy-efficient sweeping trajectories and avoid abrupt, battery-draining turns.
    \item \textbf{MVR-4: Local Moore Scan Range $r = 8$.} In 3D space, scanning up to $8$ voxels provides the ideal trade-off between local boundary awareness and computational efficiency.
    \item \textbf{MVR-5: Wind Mitigation Tuning $G \in [2.5, 3.5]$.} Flight stabilization loops must be tuned to allow moderate wind perturbations to assist with dispersion, while actively damping higher gust forces.
\end{enumerate}

% ── CONCLUSION ────────────────────────────────────────────────────────────────
\section{Conclusion \& Future Work}
\label{sec:conclusion}

This paper has presented the theoretical formalization and multi-threaded characterization of the Vel-v5.0 protocol. By extending the model to 3D continuous flight kinematics and validating it across $10{,}000{,}000$ independent trials, we have demonstrated that the protocol achieves highly stable volumetric coverage under stochastic wind disturbances. Future work will focus on deploying this 3D flight engine onto physical quadrotor hardware to validate these stability and convergence properties in real-world conditions.

% ── ACKNOWLEDGEMENTS ──────────────────────────────────────────────────────────
\section*{Acknowledgements}
The author acknowledges Gnaneswar Kumar Pullaganti and Dakshith Karthikeyan for their contributions during the early science fair phase of the Vel project. Their involvement was limited to that initial stage; they had no role in the development of Vel-v5.0, the 3D voxel flight engine, the parallel sweeps reported in this paper, or any part of this work.

The author used Gemini (Google) as an AI assistant during the preparation of this manuscript. All research concepts, experimental design, simulation implementation, data collection, and scientific conclusions are the original work of the author. The AI assistant was used for LaTeX formatting and prose editing only.

The author's ORCID is 0009-0006-2907-8815.

% ── References ────────────────────────────────────────────────────────────────
\begin{thebibliography}{99}

\bibitem{dorigo2014scholar}
M.~Dorigo, M.~Birattari, and M.~Brambilla, ``Swarm robotics,'' \emph{Scholarpedia}, vol.~9, no.~1, p. 1463, 2014.

\bibitem{pan2025}
H.~Pan, M.~Zahmatkesh, F.~Rekabi-Bana, F.~Arvin, and J.~Hu, ``T-STAR: Time-optimal swarm trajectory planning for quadrotor unmanned aerial vehicles,'' \emph{IEEE Trans. Intell. Transp. Syst.}, vol.~26, no.~8, pp. 12532--12547, 2025.

\bibitem{kernbach2013}
S.~Kernbach, Ed., ``Architectures and control of networked robotic systems,'' in \emph{Handbook of Collective Robotics}, Jenny Stanford Publishing, 2013, pp. 105--128.

\bibitem{brambilla2013review}
M.~Brambilla, E.~Ferrante, M.~Birattari, and M.~Dorigo, ``Swarm robotics: a review from the swarm engineering perspective,'' \emph{Swarm Intelligence}, vol.~7, no.~1, pp. 1--41, 2013.

\bibitem{dorigo2021past}
M.~Dorigo, G.~Theraulaz, and V.~Trianni, ``Swarm robotics: Past, present, and future [point of view],'' \emph{Proc. IEEE}, vol.~109, no.~7, pp. 1152--1165, 2021.

\bibitem{hu2020two}
J.~Hu, P.~Bhowmick, and A.~Lanzon, ``Two-layer distributed formation-containment control strategy for linear swarm systems: Algorithm and experiments,'' \emph{Int. J. Robust Nonlinear Control}, vol.~30, no.~16, pp. 6433--6453, 2020.

\bibitem{kagan2020}
E.~Kagan, Ed., \emph{Autonomous mobile robots and multi-robot systems: motion-planning, communication and swarming}, 1st ed. Hoboken, NJ: Wiley, 2020.

\bibitem{cheraghi2021}
A.~R.~Cheraghi, S.~Shahzad, and K.~Graffi, ``Past, present, and future of swarm robotics,'' \emph{arXiv preprint arXiv:2101.00671}, 2021.

\bibitem{hu2020voronoi}
J.~Hu, H.~Niu, J.~Carrasco, B.~Lennox, and F.~Arvin, ``Voronoi-based multi-robot autonomous exploration in unknown environments via deep reinforcement learning,'' \emph{IEEE Trans. Veh. Technol.}, 2020.

\bibitem{hu2020occlusion}
J.~Hu, A.~Turgut, T.~Krajnik, B.~Lennox, and F.~Arvin, ``Occlusion-based coordination protocol design for autonomous robotic shepherding tasks,'' \emph{IEEE Trans. Cogn. Dev. Syst.}, 2020.

\bibitem{alkouz2020sdaas}
B.~Alkouz, A.~Bouguettaya, and S.~Mistry, ``Swarm-based drone-as-a-service ({SDaaS}) for delivery,'' in \emph{Proc. IEEE ICWS}, 2020, pp. 441--448.

\bibitem{alkouz2020formation}
B.~Alkouz and A.~Bouguettaya, ``Formation-based selection of drone swarm services,'' in \emph{Proc. MobiQuitous}, 2020, pp. 386--394.

\bibitem{hocraffer2017}
A.~Hocraffer and C.~S.~Nam, ``A meta-analysis of human-system interfaces in unmanned aerial vehicle ({UAV}) swarm management,'' \emph{Appl. Ergon.}, vol.~58, pp. 66--80, 2017.

\bibitem{lewis2013}
M.~Lewis, ``Human interaction with multiple remote robots,'' \emph{Rev. Hum. Factors Ergon.}, vol.~9, no.~1, pp. 131--174, 2013.

\bibitem{kolling2016survey}
A.~Kolling, P.~Walker, N.~Chakraborty, K.~Sycara, and M.~Lewis, ``Human interaction with robot swarms: A survey,'' \emph{IEEE Trans. Hum.-Mach. Syst.}, vol.~46, no.~1, pp. 9--26, 2016.

\bibitem{lendon2014}
B.~Lendon, ``{U.S.} Navy could 'swarm' foes with robot boats,'' \emph{CNN}, Oct. 6, 2014.

\bibitem{madrigal2018}
A.~C.~Madrigal, ``Drone swarms are going to be terrifying and hard to stop,'' \emph{The Atlantic}, Mar. 7, 2018.

\bibitem{kushleyev2013agile}
A.~Kushleyev, D.~Mellinger, C.~Powers, and V.~Kumar, ``Towards a swarm of agile micro quadrotors,'' \emph{Autonomous Robots}, vol.~35, no.~4, pp. 287--300, 2013.

\bibitem{vasarhelyi2014}
G.~Vasarhelyi et al., ``Outdoor flocking and formation flight with autonomous aerial robots,'' in \emph{Proc. IEEE/RSJ IROS}, 2014.

\bibitem{faigl2013}
J.~Faigl, T.~Krajnik, J.~Chudoba, L.~Preucil, and M.~Saska, ``Low-cost embedded system for relative localization in robotic swarms,'' in \emph{Proc. IEEE ICRA}, 2013.

\bibitem{saska2014visual}
M.~Saska, J.~Vakula, and L.~Preucil, ``Swarms of micro aerial vehicles stabilized under a visual relative localization,'' in \emph{Proc. IEEE ICRA}, 2014.

\bibitem{saska2015}
M.~Saska, ``{MAV}-swarms: unmanned aerial vehicles stabilized along a given path using onboard relative localization,'' in \emph{Proc. ICUAS}, 2015.

\bibitem{saska2014cooperative}
M.~Saska et al., ``Autonomous deployment of swarms of micro-aerial vehicles in cooperative surveillance,'' in \emph{Proc. ICUAS}, 2014.

\bibitem{saska2014plume}
M.~Saska, J.~Langr, and L.~Preucil, ``Plume tracking by a self-stabilized group of micro aerial vehicles,'' in \emph{Modelling and Simulation for Autonomous Systems}, 2014.

\bibitem{saska2014motion}
M.~Saska, Z.~Kasl, and L.~Preucil, ``Motion planning and control of formations of micro aerial vehicles,'' in \emph{Proc. 19th World Congr. IFAC}, 2014.

\bibitem{saska2012hawk}
M.~Saska, V.~Vonasek, T.~Krajnik, and L.~Preucil, ``Coordination and navigation of heterogeneous {UAV}s-{UGV}s teams localized by a hawk-eye approach,'' in \emph{Proc. IEEE/RSJ IROS}, 2012.

\bibitem{chung2018survey}
S.-J.~Chung et al., ``A survey on aerial swarm robotics,'' \emph{IEEE Trans. Robot.}, vol.~34, no.~4, pp. 837--855, 2018.

\bibitem{saska2014mpc}
M.~Saska, V.~Vonasek, T.~Krajnik, and L.~Preucil, ``Coordination and navigation of heterogeneous {MAV}–{UGV} formations localized by a ‘hawk-eye’-like approach under a model predictive control scheme,'' \emph{Int. J. Robot. Res.}, vol.~33, no.~10, pp. 1393--1412, 2014.

\bibitem{kwon2012}
H.~Kwon and D.~J.~Pack, ``A robust mobile target localization method for cooperative unmanned aerial vehicles using sensor fusion quality,'' \emph{J. Intell. Robot. Syst.}, vol.~65, pp. 479--493, 2012.

\bibitem{itani2023}
M.~Itani, T.~Chen, T.~Yoshioka, and S.~Gollakota, ``Creating speech zones with self-distributing acoustic swarms,'' \emph{Nature Communications}, vol.~14, no.~1, p. 5684, 2023.

\bibitem{uwnews2023}
UW News, ``{UW} team's shape-changing smart speaker lets users mute different areas of a room,'' 2023.

\bibitem{harvard2014}
Harvard, ``A self-organizing thousand-robot swarm,'' Aug. 14, 2014.

\bibitem{zahugi2012}
E.~M.~H.~Zahugi, A.~M.~Shabani, and T.~V.~Prasad, ``Libot: Design of a low cost mobile robot for outdoor swarm robotics,'' in \emph{Proc. IEEE CYBER}, 2012, pp. 342--347.

\bibitem{arvin2014}
F.~Arvin, J.~C.~Murray, L.~Shi, C.~Zhang, and S.~Yue, ``Development of an autonomous micro robot for swarm robotics,'' in \emph{Proc. IEEE ICMA}, 2014, pp. 635--640.

\bibitem{beni1993}
G.~Beni and J.~Wang, ``Swarm intelligence in cellular robotic systems,'' in \emph{Proceed. NATO Advanced Workshop on Robots and Biological Systems}, Springer, 1993, pp. 703--712.

\bibitem{beni1989}
G.~Beni, ``The concept of cellular robotic system,'' in \emph{Proc. IEEE ISIC}, 1989, pp. 57--62.

\bibitem{gad2022pso}
A.~G.~Gad, ``Particle swarm optimization algorithm and its applications: A systematic review,'' \emph{Arch. Comput. Methods Eng.}, vol.~29, no.~5, pp. 2531--2561, 2022.

\bibitem{hu2021decentralized}
J.~Hu, P.~Bhowmick, I.~Jang, F.~Arvin, and A.~Lanzon, ``A decentralized cluster formation containment framework for multirobot systems,'' \emph{IEEE Trans. Robot.}, 2021.

\bibitem{sole2016}
R.~Solé et al., ``Synthetic collective intelligence,'' \emph{BioSystems}, vol.~148, pp. 47--61, 2016.

\bibitem{reynolds1987flocks}
C.~W.~Reynolds, ``Flocks, herds and schools: A distributed behavioral model,'' in \emph{Proc. SIGGRAPH '87}, 1987, pp. 25--34.

\bibitem{banks2007}
A.~Banks, J.~Vincent, and C.~Anyakoha, ``A review of particle swarm optimization. Part I: background and development,'' \emph{Nat. Comput.}, vol.~6, no.~4, pp. 467--484, 2007.

\bibitem{vicsek1995}
T.~Vicsek et al., ``Novel type of phase transition in a system of self-driven particles,'' \emph{Phys. Rev. Lett.}, vol.~75, no.~6, pp. 1226--1229, 1995.

\bibitem{buhl2006}
J.~Buhl et al., ``From disorder to order in marching locusts,'' \emph{Science}, vol.~312, no.~5778, pp. 1402--1406, 2006.

\bibitem{toner2005}
J.~Toner, Y.~Tu, and S.~Ramaswamy, ``Hydrodynamics and phases of flocks,'' \emph{Ann. Phys.}, vol.~318, no.~1, pp. 170--244, 2005.

\bibitem{reif1999}
J.~Reif and H.~Wang, ``Social potential fields: A distributed behavioral control for autonomous robots,'' \emph{Robotics and Autonomous Systems}, vol.~27, no.~3, pp. 171--194, 1999.

\bibitem{lones2014}
M.~A.~Lones, ``Metaheuristics in nature-inspired algorithms,'' in \emph{Proc. GECCO}, 2014, pp. 1419--1422.

\bibitem{sorensen2015}
K.~Sörensen, ``Metaheuristics—the metaphor exposed,'' \emph{Int. Trans. Oper. Res.}, vol.~22, no.~1, pp. 3--18, 2015.

\bibitem{glover2015}
F.~Glover and K.~Sörensen, ``Metaheuristics,'' \emph{Scholarpedia}, vol.~10, no.~4, p. 6532, 2015.

\bibitem{swan2015}
J.~Swan et al., ``A research agenda for metaheuristic standardization,'' \emph{Semantic Scholar}, 2015.

\bibitem{silberholz2019}
J.~Silberholz, B.~Golden, S.~Gupta, and X.~Wang, ``Computational comparison of metaheuristics,'' in \emph{Handbook of Metaheuristics}, Springer, 2019, pp. 581--604.

\bibitem{dorigo2006ant}
M.~Dorigo, M.~Birattari, and T.~Stutzle, ``Ant colony optimization,'' \emph{IEEE Comput. Intell. Mag.}, vol.~1, no.~4, pp. 28--39, 2006.

\bibitem{holland1975}
J.~H.~Holland, \emph{Adaptation in Natural and Artificial Systems}, Univ. Michigan Press, 1975.

\bibitem{resende2010}
M.~G.~C.~Resende and C.~C.~Ribeiro, ``Greedy randomized adaptive search procedures,'' in \emph{Handbook of Metaheuristics}, Springer, 2010, pp. 283--319.

\bibitem{kudelic2019}
R.~Kudelić and N.~Ivković, ``Ant inspired Monte Carlo algorithm for minimum feedback arc set,'' \emph{Expert Syst. Appl.}, vol.~122, pp. 108--117, 2019.

\bibitem{dorigo2004book}
M.~Dorigo and T.~Stützle, \emph{Ant Colony Optimization}, MIT Press, 2004.

\bibitem{parsopoulos2002}
K.~E.~Parsopoulos and M.~N.~Vrahatis, ``Recent approaches to global optimization problems through particle swarm optimization,'' \emph{Nat. Comput.}, vol.~1, pp. 235--306, 2002.

\bibitem{clerc2006}
M.~Clerc, \emph{Particle Swarm Optimization}, ISTE, 2006.

\bibitem{rosenberg2015human}
L.~Rosenberg, ``Human swarms, a real-time method for collective intelligence,'' in \emph{Proc. European Conf. Artificial Life}, 2015, pp. 658--659.

\bibitem{metcalf2019}
L.~Metcalf, D.~A.~Askay, and L.~B.~Rosenberg, ``Keeping humans in the loop: Pooling knowledge through artificial swarm intelligence to improve business decision making,'' \emph{Calif. Manage. Rev.}, vol.~61, no.~4, pp. 84--109, 2019.

\bibitem{rosenberg2018diagnostic}
L.~Rosenberg et al., ``Artificial swarm intelligence employed to amplify diagnostic accuracy in radiology,'' in \emph{Proc. IEEE IEMCON}, 2018, pp. 1186--1191.

\bibitem{alrifaie2012SDS}
M.~M.~al-Rifaie and A.~Aber, ``Identifying metastasis in bone scans with stochastic diffusion search,'' in \emph{Proc. IEEE ITME}, 2012, pp. 519--523.

\bibitem{sun2020}
W.~Sun et al., ``A survey of using swarm intelligence algorithms in {IoT},'' \emph{Sensors}, vol.~20, no.~5, p. 1420, 2020.

\bibitem{abualigah2023}
L.~Abualigah, D.~Falcone, and A.~Forestiero, ``Swarm intelligence to face {IoT} challenges,'' \emph{Comput. Intell. Neurosci.}, 2023.

\bibitem{martens2011}
D.~Martens, B.~Baesens, and T.~Fawcett, ``Editorial survey: Swarm intelligence for data mining,'' \emph{Mach. Learn.}, vol.~82, no.~1, pp. 1--42, 2011.

\bibitem{thrun2021}
M.~Thrun and A.~Ultsch, ``Swarm intelligence for self-organized clustering,'' \emph{Artificial Intelligence}, vol.~290, p. 103237, 2021.

\bibitem{miller2010}
P.~Miller, \emph{The Smart Swarm}, New York: Avery, 2010.

\bibitem{shah2023}
R.~Shah et al., ``Utilizing a digital swarm intelligence platform to improve consensus among radiologists and exploring its applications,'' \emph{J. Digit. Imaging}, vol.~36, no.~2, pp. 401--413, 2023.

\bibitem{mammen2009}
S.~von Mammen and C.~Jacob, ``The evolution of swarm grammars -- growing trees, crafting art and bottom-up design,'' \emph{IEEE Comput. Intell. Mag.}, vol.~4, no.~3, pp. 10--19, 2009.

\bibitem{rifaie2012creative}
M.~M.~al-Rifaie, J.~M.~Bishop, and S.~Caines, ``Creativity and autonomy in swarm intelligence systems,'' \emph{Cogn. Comput.}, vol.~4, no.~3, pp. 320--331, 2012.

\bibitem{deleuze2004}
G.~Deleuze, F.~Guattari, and B.~Massumi, \emph{A Thousand Plateaus}, Univ. Minnesota Press, 2004.

\bibitem{rifaie2013sketches}
M.~M.~Al-Rifaie and J.~M.~Bishop, ``Swarmic sketches and attention mechanism,'' in \emph{Evolutionary and Biologically Inspired Music, Sound, Art and Design}, Springer, 2013, pp. 85--96.

\end{thebibliography}

% --- Author Biography Section ---
\section*{Author Biography}
% The [-2ex] reduces the gap between the last section and the biography
\vspace{-7ex} 

\begin{IEEEbiographynophoto}{Dakshith Rajkumar}
is an independent researcher specializing in decentralized swarm intelligence and autonomous robotic navigation. His work focuses on the development of the Vel protocol series, specifically exploring Map-Aware Gradient Repulsion (MAGR) and spatio-temporal efficiency in multi-agent systems. A member of the National Junior Honor Society (NJHS) at Fowler Middle School, part of the Frisco Independent School District, he maintains a rigorous commitment to academic excellence and emergent engineering logic.
\end{IEEEbiographynophoto}

\vfill % Pushes the bio to the bottom or keeps spacing consistent with IEEE style

\end{document}